% Table T18 -- All Features versus Selected Features
% Experiment: EXP-F1
% Objective: O5 (feature selection evidence)
% Source: outputs/09_ablation/feature_ablation_per_fold.csv; outputs/03_features/all_features_matrix.parquet; outputs/01_dataset_audit/subject_split_map.csv; outputs/03_features/selected_feature_subset.csv
% Generated by PV-MEPCG / PulseVision at 2026-09-10T07:31:27.078675+00:00
\begin{table}[htbp]
  \centering
  \caption{A7 (all 138 features) against A8 (the SO-04 selected subset) on the identical DA-07 repeated 5x5 subject-grouped fold map, identical seed and identical pipeline. A positive delta means the subset scored higher than the full representation. Deltas are PAIRED per fold -- the two arms ran the same folds -- so the SD is the spread of the difference itself, and 'folds won' counts the folds on which the subset beat all 138.}
  \label{tab:t18}
  \begin{tabular}{lrrrrrrrrrrrrrrrrrrrrrrrrrrrrrrrrr}
    \toprule
    Model & Features (all) & Features (selected) & Reduction & sensitivity (all 138) & sensitivity (selected) & sensitivity Δ & sensitivity Δ SD & sensitivity folds won & specificity (all 138) & specificity (selected) & specificity Δ & specificity Δ SD & specificity folds won & f1 (all 138) & f1 (selected) & f1 Δ & f1 Δ SD & f1 folds won & balanced accuracy (all 138) & balanced accuracy (selected) & balanced accuracy Δ & balanced accuracy Δ SD & balanced accuracy folds won & roc auc (all 138) & roc auc (selected) & roc auc Δ & roc auc Δ SD & roc auc folds won & accuracy (all 138) & accuracy (selected) & accuracy Δ & accuracy Δ SD & accuracy folds won \\
    \midrule
    M1 & 138 & 20 & 0.855 & 0.848 & 0.873 & 0.025 & 0.032 & 19 & 0.838 & 0.802 & -0.036 & 0.020 & 1 & 0.687 & 0.663 & -0.023 & 0.023 & 3 & 0.843 & 0.837 & -0.005 & 0.016 & 7 & 0.922 & 0.911 & -0.011 & 0.009 & 2 & 0.840 & 0.816 & -0.024 & 0.016 & 3 \\
    M5 & 138 & 20 & 0.855 & 0.798 & 0.804 & 0.006 & 0.023 & 13 & 0.902 & 0.896 & -0.005 & 0.011 & 6 & 0.734 & 0.730 & -0.004 & 0.022 & 13 & 0.850 & 0.850 & 0.000 & 0.014 & 14 & 0.936 & 0.936 & 0.000 & 0.006 & 10 & 0.880 & 0.877 & -0.003 & 0.011 & 8 \\
    M4 & 138 & 20 & 0.855 & 0.748 & 0.787 & 0.040 & 0.022 & 24 & 0.937 & 0.924 & -0.013 & 0.009 & 2 & 0.751 & 0.758 & 0.007 & 0.020 & 17 & 0.842 & 0.856 & 0.014 & 0.012 & 23 & 0.942 & 0.945 & 0.003 & 0.004 & 21 & 0.898 & 0.896 & -0.002 & 0.009 & 10 \\
    M8 & 138 & 20 & 0.855 & 0.756 & 0.772 & 0.017 & 0.024 & 19 & 0.933 & 0.926 & -0.007 & 0.009 & 5 & 0.750 & 0.751 & 0.001 & 0.020 & 13 & 0.844 & 0.849 & 0.005 & 0.014 & 16 & 0.941 & 0.944 & 0.003 & 0.006 & 20 & 0.896 & 0.895 & -0.002 & 0.009 & 9 \\
    M3 & 138 & 20 & 0.855 & 0.670 & 0.642 & -0.028 & 0.043 & 4 & 0.930 & 0.934 & 0.004 & 0.007 & 18 & 0.691 & 0.676 & -0.015 & 0.029 & 7 & 0.800 & 0.788 & -0.012 & 0.021 & 6 & 0.924 & 0.928 & 0.004 & 0.006 & 20 & 0.877 & 0.874 & -0.003 & 0.009 & 8 \\
    \bottomrule
  \end{tabular}
  \par\footnotesize The subset was chosen ONCE, by FE-12, under a one-standard-error performance guard, and is read from disk here. It was not re-selected per fold inside this comparison; a subset re-chosen with the test folds visible would beat all 138 by construction.
  \par\footnotesize A reduction figure is not a speed claim. The cost of computing the 20 retained features is not 20/138 of the cost of computing all 138 -- sample entropy dominates extraction and lives in the time family. T24-T26 carry the measured timings.
  \par\footnotesize PV-MEPCG / PulseVision is an academic screening prototype, not a diagnostic tool.
\end{table}
